\section{Introduction}
\label{sec:intro}

Real-time optimisation (RTO) recomputes the operating point of a process to maximise an
economic objective as conditions drift. Its value comes from pushing the process against
its constraints; a distillation column earns more as it drives a product composition to
the edge of specification. The constraints that matter most, however, are often on quality
variables that are not measured online. A laboratory assay or an inferential estimate
stands in for a true composition that the controller never sees in real time. To operate
safely against such a constraint, an RTO must subtract a back-off from the specification: a
margin that absorbs the gap between the model's nominal prediction and the realised quality
under an unmeasured disturbance. Sizing that margin is the central difficulty of RTO under
uncertainty. Set it too small and the process violates specification whenever the
disturbance is adverse. Set it too large and the economic benefit of the optimisation is
given away.

Conformal prediction is the natural tool for the job. It turns any predictor into an
interval with a finite-sample, distribution-free coverage guarantee, assuming only that the
data are exchangeable, with no Gaussian residuals and no parametric disturbance law. Taking
the upper edge of a conformal interval as the back-off promises exactly what a chance
constraint asks for, namely coverage of the realised quality at the nominal level, learned
from data, without distributional assumptions. This paper shows that the promise does not
survive contact with the optimiser, and it supplies a back-off construction that restores it.

An RTO that optimises against a marginally valid conformal back-off is unsafe. A conformal
margin guarantees coverage on average over the calibration distribution of operating points,
but the optimiser does not sample that distribution. It deliberately seeks the boundary of
the feasible region, where the margin is tightest relative to the local risk and therefore
where coverage is worst. The argmax is a selection operator, and it selects for
under-coverage. On a rigorous binary-distillation twin (Figure~\ref{fig:rto}), a marginally
valid back-off, including a locally adaptive heteroscedastic one that widens where the data
are noisier, yields a realised violation rate of roughly five times the nominal level across
the entire disturbance range; the failure is structural rather than a consequence of a large
disturbance, and adaptivity in the width of the margin does not cure it. The mechanism is an
optimisation-induced loss of marginal coverage, a constraint-space analogue of the winner's
curse in post-selection inference. The conformal-predictive-programming literature
\citep{zhao2024cpp} observed in the abstract that a decision optimised against sampled
constraints is no longer independent of them; our concern here is the concrete consequence
for process RTO, a safety hazard of quantified severity, and, more to the point, the back-off
construction that removes it.

What fixes it is conditional validity. A back-off built from conformalised quantile
regression \citep{romano2019} targets the conditional constraint quantile at each operating
point rather than a calibration average, so there is no systematic gap between margin and
risk for the optimiser to exploit. The residual finite-sample error is closed by an
a-posteriori calibration at the selected setpoint, in the manner of conformal predictive
programming. On the twin, this conditional method returns the realised violation to the
level of an oracle that knows the true conditional quantile, to within the Monte-Carlo
estimation band, at near-oracle profit over the operationally realistic disturbance range.
The oracle margin, available here because the twin furnishes the plant response, holds the
violation at the nominal level by construction. That it does so confirms that the failure
of the marginal methods lies in the back-off, not in the chance-constrained formulation.

The method is not unconditionally safe, and we map where it holds. Sweeping the disturbance
magnitude produces a regime map. The conditional method tracks the oracle up to roughly
three times the realistic disturbance. A constant margin over-violates and then becomes
infeasible. The a-posteriori inflation factor climbs from near unity to nearly sixfold
before the conditional estimate, too, runs out of calibration data and can no longer
certify the target. That inflation factor is an interpretable diagnostic, since the operator
sees the safety margin being stretched before it breaks. Finally, we are explicit that the
contribution is calibrated safety, not profit. At a well-controlled feed the constraint is
barely active and every method earns within half a percent of the deterministic optimum, so
the value of the conditional method is the violation it removes, not any profit it adds.

The contributions are as follows.
\begin{itemize}
\item \textbf{R1.} We propose a deployable, distribution-free constraint back-off for
  chance-constrained RTO, built from conformalised quantile regression with an a-posteriori
  calibration at the selected setpoint, and show that it holds the realised violation at the
  oracle level, within the Monte-Carlo band, at near-oracle profit over the operationally
  realistic disturbance range (Section~\ref{sec:results-fix}).
\item \textbf{R2.} We identify and quantify the deployment hazard that motivates it. A
  marginally valid back-off, fixed or locally adaptive, realises roughly five times the
  nominal violation under optimisation, because the optimiser settles where the margin
  under-covers the conditional risk; an oracle conditional margin holds the target, which
  establishes that the failure lies in the back-off and not in the chance-constrained
  formulation (Sections~\ref{sec:results-selection} and \ref{sec:results-oracle}).
\item \textbf{R3.} We characterise the disturbance-magnitude regime over which each margin
  controls violation, where a constant margin becomes infeasible, and where the conditional
  method itself is data-starved, and we show that the a-posteriori inflation factor serves as
  an interpretable data-adequacy diagnostic, visible to the operator before the margin breaks
  (Sections~\ref{sec:results-boundary} and \ref{sec:results-profit}).
\end{itemize}

Section~\ref{sec:related} places the work against RTO under uncertainty, conformal
prediction, and conformal optimisation. Section~\ref{sec:methods} formalises the
chance-constrained RTO, the four back-offs, the selection-effect argument, and the
a-posteriori step. Section~\ref{sec:case} describes the twin and the experimental design,
Section~\ref{sec:results} reports the regime map, and Section~\ref{sec:discussion} draws out
the general design rule, the limitations, and the bridge to a closed-loop extension with an
online soft sensor.
